\chapter*{How to Read This Companion}
\phantomsection\addcontentsline{toc}{chapter}{How to Read This Companion}
\markboth{How to Read This Companion}{}
\label{comp:guide}

\section*{What this companion is}
\phantomsection\addcontentsline{toc}{section}{What this companion is}

The theory lives in three technical documents: the main manuscript, cited here as \textsc{mt}, which
develops everything at one base point $c_*$; Supplement~G, the multivariate-Gaussian realization; and
Supplement~B, which holds base transport, infinite networks, and the lattice cluster, all deferred
until the single-point theory is settled. Together they run to more than six hundred pages. They are
written for correctness under adversarial checking, not for a first reading, and they are the only
authority. This companion proves nothing. It explains what the technical documents say, in the
language of a physicist, and it says plainly what kind of knowledge each main result is and whether
the theory delivers what its own words promise.

When the technical documents prove something, the companion states it and cites it. A citation such
as \MT{thm:elbo-extended-gap} refers to the main manuscript, and statements of the supplements carry
their own letters (Proposition~G4.2, Theorem~B6.1). In the PDF these references are links that open
the technical document at the cited statement, provided the technical PDFs sit in the neighboring
folder where they are built. The companion's main part mirrors the sixteen chapters of the main
manuscript one to one, including the latent-link chapter (\cref{comp:lie-links}). Three chapters cover
Supplement~G (\cref{comp:sg1,comp:sg2,comp:gaussian-latent-links}), and \cref{comp:sb} covers Supplement~B. The
companion has no numbered statements of its own, so every theorem, proposition, or lemma number you
meet here points into a technical document.

Each companion chapter follows one template: the question the technical chapter answers and its
answer in plain language; the physical picture, with the nearest physics analogue graded as exact,
structural, or loose; a decoder for the chapter's notation and jargon; the key formulas explained
term by term with their limiting cases; a scorecard of the main results; a box that sorts the
chapter's hypotheses into those that matter and those that are hygiene; the open ends; and a
verdict. Where a proof would help intuition, the companion gives a sentence or two of informal
reason and labels it as such.

Highlighted equation boxes mark the central constructions and balances of the dynamic-link model.
Their titles identify the role of each formula; the surrounding text states its hypotheses and
points to the technical result. The Proca-type equation in \cref{comp:lie-links} is a linearized
Euclidean stationarity equation across inference episodes, whereas the recognition flow is first order.

\section*{Two questions about every result}
\phantomsection\addcontentsline{toc}{section}{Two questions about every result}

The first question is what kind of knowledge a result is. Nearly every theorem-type statement in the
main manuscript carries the tag \textsc{established}, meaning it is proved there. The tag says
nothing about how general the result is, how much it depends on a special instance, or whether it
rests on an assumption that might fail in the cases the program cares about. The companion adds that
information with one badge per result.

\begin{sciencetable}{Status badges used in this companion.}
\label{comp:tab-badges}
\begin{tabularx}{\textwidth}{>{\raggedright\arraybackslash}p{0.13\textwidth}>{\raggedright\arraybackslash}X}
\toprule
\textbf{Badge} & \textbf{Meaning} \\
\midrule
\SOLID & Proved in the technical documents, or a checked textbook result, at the generality the theory relies on. Established negative results (refutations and no-go statements) are solid. When the mathematics is textbook and the novelty is only the typing, the companion says ``standard in substance''. \\
\TOY & Proved, but only on a named small or special instance: the four-cycle, a pair, three binary agents, a scalar or abelian case, a particular constant schedule. The mechanism is not established beyond that instance. \\
\COND & Proved under a hypothesis that could fail in the models the program cares about and is not shown to hold there. The companion names the hypothesis. \\
\NUM & Supported only by a named computation. \\
\OPEN & An open problem or a conjecture. \\
\VAGUE & A reading, an interpretation, or a programmatic statement with no precise mathematical statement behind it, or a construction whose meaning is not pinned down. \\
\CHOICE & A definition or standing hypothesis that is a substantive modeling commitment, the place where physics enters, as opposed to notation. \\
\bottomrule
\end{tabularx}
\end{sciencetable}

The second question is whether the theory does what it says. Each companion chapter ends with a
verdict on the technical chapter's own claims, taken from its opening paragraph and from the reading
map of the introduction, measured against what its statements actually prove.

\begin{sciencetable}{Delivery verdicts.}
\label{comp:tab-verdicts}
\begin{tabularx}{\textwidth}{>{\raggedright\arraybackslash}p{0.27\textwidth}>{\raggedright\arraybackslash}X}
\toprule
\textbf{Verdict} & \textbf{Meaning} \\
\midrule
\DELIVERED & The words match the mathematics. \\
\NARROWER & The claim is proved, but in a setting narrower than the headline wording suggests: a supplied hierarchy, a single small network, a particular family. \\
\PARTIAL & Part of the claim is proved and part is open. \\
\NOTDELIVERED & The claim is posed but not proved; the technical document usually says so itself. \\
\bottomrule
\end{tabularx}
\end{sciencetable}

A result that is correct and standard is not a defect, and the companion says ``standard in
substance'' as information, not as criticism: a physicist gains from knowing that an identity is the
Gibbs variational principle in new notation, because then the physicist already knows what it can
and cannot do.

\section*{Reading the technical documents: the house style decoded}
\phantomsection\addcontentsline{toc}{section}{Reading the technical documents: the house style decoded}

The technical documents are written under an authoring specification whose rules explain most of
what makes them hard to read. Knowing the rules lets you skip what is bookkeeping and slow down where
a qualifier carries content.

\paragraph{Status tags and scope paragraphs.}
Every lemma, proposition, theorem, corollary, and claim-bearing remark carries exactly one tag:
\textsc{established} (proved in the document), \textsc{standard} (a known result, cited with a
locator), or \textsc{numerical} (supported only by a named computation). Definitions, standing
hypotheses, conjectures, and open problems carry their status in the heading. After its opening, a
chapter may add one scope paragraph, tagged \textsc{not-claimed}, listing what it declines to claim;
in practice every chapter has one. Declining is
not refuting: a declined statement may be true, and a refuted statement appears instead as an
established negative result. The scope paragraphs are the fastest way to find out what a chapter does
not do, and the companion's verdicts lean on them.

\paragraph{Defaults, and why the prose is full of qualifiers.}
The introduction of the main manuscript fixes defaults once for the whole document
(\MT{sec:intro-conventions}): measurable spaces are standard Borel, laws and kernels are normalized
and measurable, a coarse channel is a Markov kernel that does not read the recognition law and
leaves observations alone, admitted observations have finite positive evidence, structural data are
parameters unless a law for them is declared, and the interaction multigraph contains the reversal
closure of the union of the source-prior supports. A receiving row's support is complete unless
masked and need not be symmetric. In the latent-link model the multigraph and frames remain
structural, while Lie-algebra generators are inferred under a prior
(\MT{def:link-latent-model}). A statement mentions only departures from these defaults, so a
hypothesis that appears inside a statement is a departure and worth reading. Three other sources of
qualifiers are worth knowing. The documents work with arbitrary laws on arbitrary spaces rather than
Gaussians, which forces extended-real values (a relative entropy may be $+\infty$) and almost-sure
qualifiers. They keep apart objects that coincide in the Gaussian or flat cases but not in general:
three description levels (sample, law point, configuration), three transports (frame comparison,
graph link, base parallel transport), three kinds of agreement between agents, three coarse
operations (Markov pushforward, energy precomposition, Galerkin restriction), and five attention objects
(the source prior, the generative and the recognition attention rows, the receiver occupancy, and the
event law).
Keeping them apart is deliberate: each pair can differ outside the Gaussian or flat case, and a
theorem about one does not transfer to the other. And the scope rules described next constrain what
may be called formation, selection, or recursion.

\paragraph{The Curie rule (R13): what may be called selection by learning.}
Curie's principle says that symmetric causes cannot produce less symmetric effects when the effect
is uniquely determined by the cause. The rule applies it to learning. Let $H$ be the symmetries shared
by the generative model with its records, the parameterization, and the kinetics, including the
kinetic metric and the proposal mechanism ($H$ may be a group of graph automorphisms or a residual
gauge group). A claim that learning creates or selects a structure that $H$ does not fix must name
the $H$-breaking datum it amplifies (the initial law, an initial parameter, the model, or the
records) and must compare with the frozen-learning arm on the same data. A claim that learning
selects a structure that $H$ does fix must name the symmetry it restores, with the frozen arm at the
non-invariant initial value as its null. A claim of selection with no breaking datum must name the
hypothesis its mechanism violates and declare the mechanism as part of the model: a rule that reads
realized paths, records, or a finite ensemble; selection among non-unique optimizers, critical
points, or branches; non-uniqueness at masked entries, which amounts to edge creation; or a
non-equivariant metric, mobility, or proposal mechanism. This is substantive, not bookkeeping: an
equivariant law-level learning rule whose coupled law and parameter dynamics is unique cannot break a
symmetry that its initial law and initial parameter share. In physics terms (the companion's
reading, not part of the rule), spontaneous symmetry breaking needs non-uniqueness, which in
statistical mechanics arrives with the thermodynamic limit. The four-cycle results of
\cref{comp:ch14} are written under this rule, which is why they speak of amplifying an initial
asymmetry rather than of spontaneous selection.

\paragraph{Recursion versus a supplied-scale sequence (R14).}
A two-stage claim declares a handoff time, the first- and second-stage events with their blocks, one
statistic of the handoff state, and a class of handoffs fixed in advance that also contains, with
probability bounded below, handoffs of histories in which the first stage did not form. If the
second-stage event occurs with probability tending to one uniformly over the handoff class, the
sequence is a supplied-scale sequence, not recursion, whatever differences appear earlier. Otherwise
the second coalescence counts as recursive only if, after first-stage formation, the second-stage
event has probability bounded below, its block is a union of at least two first-stage blocks, and
the declared statistic exceeds its value after a counterfactual handoff by a margin bounded below. A
margin in the other direction, where the first stage delays or suppresses the second, is recorded as
first-stage dependence, not recursion. In physics terms (again the companion's reading), a
supplied-scale sequence is a hierarchy built into the interaction strengths through which the
dynamics relaxes in order, and the rule separates a genuine cascade from the successive freezing of
modes whose stiffnesses were prescribed. It matters most in \cref{comp:ch13}.

\paragraph{Coherence episodes (R15): what counts as formation.}
A block counts as having formed only if two clauses hold. The restoring clause names the interaction
that holds the block coherent against thermal noise at the stated resolution, with restoring rate
times squared resolution divided by temperature tending to infinity, and the rate is evaluated on the
actual law of the assignment sheets or on the landscape marginalized over sources, not on a sheet
held fixed by hand. The horizon clause requires the
coherence to be a dynamical event rather than a property of equilibrium: either the observation
horizon is short compared with the stationary recurrence time of the event, or the event is more
probable than it is in the stationary state by a margin. Coherence that is restored but merely
reflects equilibrium (a strongly bound pair in thermal equilibrium is always coherent) is called
restored equilibrium coherence; coherence with no restoring interaction is coincidental. In physics
terms, the rule separates ordering produced by the dynamics from equilibrium order and from
fluctuations.

\paragraph{Settling conditions (R16).}
Every open problem and conjecture states a class of models and two outcomes, a proof or a
counterexample within that class, and neither outcome may already be met by an existing statement or
by a trivial construction. This is why the open problems read like contracts. It keeps open items
falsifiable, at the price of posing them in narrower classes than the questions that motivate them.

\paragraph{No Wilsonian vocabulary in the main manuscript.}
The main manuscript constructs exact finite-step coarse-graining with no rescaling, so it does not
use the words relevant, irrelevant, fixed point, or beta function except in one pointer to
Supplement~B and when describing other authors' work. The asymptotic apparatus (rescalings, spectra, fixed points, universality) lives in
Supplement~B. When the companion uses renormalization-group language, it says which of the two it
means.

\paragraph{Labels and names.}
Labels are frozen when a statement changes type, so a label beginning \texttt{thm:} may belong to a
corollary and \texttt{conj:phil-contrast-formation} is a theorem. Trust the printed heading. Words
such as agent, meta-agent, parent, and block are naming decisions: the documents make no claim of
agency, and population laws are not agents.

\section*{A notation survival kit}
\phantomsection\addcontentsline{toc}{section}{A notation survival kit}

The notation table of the main manuscript (\MT{app:notation-contract}) is complete but long. The
following symbols carry most of the load; the glossary at the end of the companion
(\hyperref[comp:glossary]{Glossary}) covers the jargon.

\begin{decoder}
$c_*$ & The single base point (context) at which everything happens. & A zero-dimensional base; one spacetime point. \\
$G$, $\mathscr P_G$ & Structure group and principal bundle; at $c_*$ the fiber is a $G$-torsor (a copy of $G$ with no preferred identity). & Gauge group; the space of frames at a point. \\
$\rho_b$, $\rho_m$ & Representations of $G$ on the belief and model sample spaces; they may differ. & Matter representations. \\
$q_i^b$, $q_i^m$ & Agent $i$'s belief law and model law (probability laws on the fibers). & A mixed state of a matter field at site $i$. \\
$u_i$ & Agent $i$'s frame, one for both channels at $c_*$ (a gauge choice). & A local gauge fixing at site $i$. \\
$\mathfrak I$, $e$, $\bar e$ & Interaction multigraph, an oriented edge copy, its reversal. & The lattice, with multiple links allowed. \\
$\Theta^x_e$ & Independent group-valued link on edge copy $e$ for channel $x\in\{b,m\}$, with $\Theta_{\bar e}=\Theta_e^{-1}$. & Link variable $U_\ell$ of lattice gauge theory. \\
$\mathring\omega^x_e$, $R$ & Generator of the reference twist, $\Theta^x_e=U_i\exp(\mathring\omega^x_e)U_j^{-1}$; recognition law of the generators. & Lie-algebra field and its variational distribution. \\
$\mathring\alpha$, $\mathring\Sigma$ & Anchor and covariance of the Gaussian base prior; generative parameters. & Preferred field and bare restoring scale. \\
$\mathscr S$, $\mathbf m$, $n$ & Loop action, history increment precision, and learning-episode index. & Spatial loop penalty, Euclidean history mass, and an inference index. \\
$\mathrm{Hol}^x(p)$ & Holonomy: product of links around a closed walk $p$. & Wilson loop before the trace. \\
$K_a$, $\mathfrak A$ & A normalized record kernel (an interaction or observation factor) and the set of them. & A factor of the Boltzmann weight; a local term of the action. \\
$\mathbb P_{\theta,V}$ & The one normalized joint generative law of all agents and records. & The Gibbs measure $e^{-S}/Z$. \\
$\mathbb Q_{V,o,X}$ & The one recognition law of the whole population (a coupling of the agents' local laws). & The variational trial state. \\
$\boldsymbol\Pi_{\theta,V,o,X}$ & The exact posterior given observation $o$ and structure $X$. & The exact Gibbs state conditioned on the data. \\
$\Lelbo_V$, $\Fenergy_V=-\Lelbo_V$ & Evidence lower bound and variational free energy of the population. & Minus the variational free energy; the Gibbs--Bogoliubov functional. \\
$\pi_{ij}$, $\beta_{ij}$, $\tau_i$ & Source prior, attention weight of receiver $i$ on source $j$, attention temperature. & Prior weight, Boltzmann weight over sources, temperature. \\
$S^Q_e$, $S^K_e$, $\mathsf L_e$ & Query and key comparison maps and the composite communication map of an edge. & None standard; linear maps between internal spaces. \\
$C_A$, $\ell$ & Coarse channel that makes parent $A$ from its child block; coarse-graining depth. & Blocking kernel; step index of a renormalization without rescaling. \\
$\mathrm T$, $\varepsilon$ & Temperature and small parameters of the formation families (\cref{comp:ch13,comp:ch14}). & Temperature; small expansion parameters. \\
\end{decoder}

\section*{Reading paths}
\phantomsection\addcontentsline{toc}{section}{Reading paths}

For a first pass, read \hyperref[comp:glance]{The Theory at a Glance} and the verdict box at the end of each chapter. For the
physical core of the theory, read \cref{comp:ch02,comp:ch05,comp:ch06,comp:lie-links,comp:ch09,comp:ch11,comp:ch13,comp:ch14}
in that order. For what the theory says about Gaussian networks with $\GL(K)$ links at large size,
read \cref{comp:ch02,comp:lie-links,comp:ch10,comp:sg1,comp:sg2,comp:gaussian-latent-links}. For the state of the renormalization program at one
base point, read \cref{comp:ch09,comp:ch11,comp:ch12,comp:sb}. The open ends of the whole theory are
collected in \hyperref[comp:open-ends]{Open Ends of the Theory}.
